\subsection{A prime-weight invariant of the single-disk localization}

The restriction $N\geq2$ in the definition of
$\mathcal E^{[d]}$ is retained throughout this subsection.
The following calculation determines the proposed localization
comparison for this actual category.

\begin{theorem}[A-S09-28: Proved; the single-disk inclusion is not faithful after localization]
\label{AS0928}
For every $d\geq0$, the all-arrow localization of
$\mathcal E^{[d]}$ has a nonidentity automorphism.
The functor induced by the actual full inclusion
\[
 \mathcal E^{[d]}\longrightarrow\mathcal A_{\mathbb Z}
 \tag{S09.147}
\]
is not faithful on all-arrow localizations.
More precisely, an explicit object $E\in\mathcal E^{[d]}$
and an explicit strict endomorphism $u:E\to E$ have
$[u]\ne1$ in the first localization and $[u]=1$ in the
second. In fact the same endomorphism already becomes the identity
in the localization of $\mathcal C^{[d]}$, by the explicit
doubled-skeleton diagram below. The same assertions hold
with support $q\leq w$, termwise integral flatness, or both.

Thus the common-source equivalence of A-S09-27 cannot be
promoted to an equivalence of these particular localizations.
This result leaves the terminal-localization gate
$\mathcal C^{[d]}$ of A-S09-24 unchanged.
\end{theorem}

\begin{proof}
We first specify an integral functor on first skeletons.
For $P(N)_1^0$, let $J$ be its positive-weight ideal.
Quotient $J$ by the subgroup generated by its products,
all $\gamma_r(z)$ with $z\in J$ and $r\geq2$, and the
positive-weight coefficient ideal. Denote the quotient by
$Q(P(N))$. The free PD monomial presentation in (SK.1)
gives the actual identification
\[
             Q(P(N))=\bigoplus_{w\geq N}\mathbb Z\,\overline{x}_w.
 \tag{S09.148}
\]
Indeed a monomial with a positive coefficient and a variable
is a product; a monomial with several positive variable
factors is a product; and a single divided power of order
at least two is one of the listed relations. Pure coefficients
are killed. These operations leave exactly the linear
monomials $x_w$, which are independent in the free presentation.
Conversely the map taking a PD polynomial to its linear
variable part kills every listed relation, including
$\gamma_r(z)$ for arbitrary $z\in J$, by the PD addition
formula. This proves both directions of (S09.148).
Every strict coefficient-preserving PD map preserves the
listed subgroups, so this is an actual functor.

Let $f:P(M)\to P(N)$ be any strict first-skeleton map.
For each weight $w\geq\max(M,N)$, write
\[
                    Q(f)(\overline{x}_w)
                         =a_w(f)\overline{x}_w,
                    \qquad a_w(f)\in\mathbb Z.
 \tag{S09.149}
\]
In particular these scalars multiply under composition.
If $p\geq\max(M,N)$ is prime and $M,N\geq2$, then
\[
                            a_p(f)\equiv1\pmod p.
 \tag{S09.150}
\]
Here is the complete endpoint calculation. The condition
$d_1f(x_p)=0$ removes its pure coefficient term, since
$d_1$ kills every variable and is the identity on coefficients.
Expand the remaining terms in the free PD monomial basis
\[
       T^{[k]}\prod_{v\geq N}\gamma_{r_v}(x_v),
               \qquad k+\sum_v r_vv=p,
               \qquad\sum_v r_v>0.
\]
The term $x_p$ has zeroth face $T^{[p]}$.
For any other displayed monomial all nonzero $v$ are
less than $p$, all $r_v$ are less than $p$ because
$v\geq2$, and $k<p$. Its zeroth face is exactly
\[
        \frac{p!}{k!\prod_v r_v!\,(v!)^{r_v}}\,T^{[p]}.
 \tag{S09.151}
\]
The coefficient is an integer by the divided-power
multiplication formulas. Its denominator is prime to $p$,
whereas $p!$ is divisible by $p$. Thus that integer is
divisible by $p$. Since $d_0f(x_p)=T^{[p]}$, reduction
of this actual endpoint equality proves (S09.150).
In particular $a_p(f)\ne0$.
Using prime weights is essential in this calculation;
no analogous claim for all prime powers is being made.

Let
\[
 \mathcal G=
 \left(\prod_{p\ \mathrm{odd\ prime}}\mathbb Q^\times\right)
 \Big/
 \left\{(c_p):c_p=1\text{ for all but finitely many }p\right\}.
 \tag{S09.152}
\]
This is an abelian group. For each object of
$\mathcal E^{[d]}$ use a first-skeleton identification
with its $P(N)$, and for each morphism use the germ of
its scalars (S09.149) in all prime weights above the
two cutoffs. Set the finitely many omitted components
equal to one. Equation (S09.150) makes every retained
component a member of $\mathbb Q^\times$, and functoriality
of $Q$ proves that this gives a functor
\[
 \chi:\mathcal E^{[d]}\longrightarrow B\mathcal G,
 \tag{S09.153}
\]
where $B\mathcal G$ is the one-object groupoid with
morphism group $\mathcal G$.

There is no dependence on a hidden choice of skeleton
orientation. First, the first weight in which (S09.148)
is nonzero determines $N$. An isomorphism of two
presentations with that same $N$ induces an automorphism
of the rank-one integral group in each weight, hence a
scalar $\pm1$. In every sufficiently large odd prime
weight (S09.150) forces that scalar to be $+1$.
Changing the identifications therefore changes only
finitely many of the prime components. The germ in
(S09.152) is unchanged. The functor sends every arrow
to an isomorphism and consequently extends to the
all-arrow localization by sending a formal inverse
to the inverse element of $\mathcal G$.

We now give the object and endomorphism detected by
this functor. Choose
\[
        N\geq\max\{2,\lambda_G(d+2),\mu_G(d+2)\},
\]
where $G$ is the standard integral theory. In
$G_1^0=S\otimes\Gamma(x)$ the faces of $x$ are
$0,T$. Define the strict truncated map
$p:P(N)\to(T_NG)_{\leq1}$ by
$x_w\mapsto\gamma_w(x)$ and $y_w\mapsto d\gamma_w(x)$.
Its values have weight at least $N$, so they are actual
elements of the unchanged high-weight terms of $T_NG$.
It is coefficient preserving and has the prescribed
endpoints. Use precisely the A-S09-27 replacement
\[
 E=T_NG\mathop{\times}_{C((T_NG)_{\leq1})}C(P(N)),
                       \qquad \pi:E\longrightarrow T_NG.
\]
That theorem proves $E\in\mathcal E^{[d]}$, including
the support and flatness assertions for this object.

Define an actual strict endomorphism $\psi$ of $P(N)$
by the free PD generator assignments
\[
 \begin{aligned}
 \psi(x_w)&=x_w &&(N\leq w<2N),\\
 \psi(x_w)&=\left(1+\binom wN\right)x_w-x_Nx_{w-N}
                                             &&(w\geq2N),\\
 \psi(y_w)&=d\psi(x_w),\qquad \psi|_S=\mathrm{id}.
 \end{aligned}
 \tag{S09.154}
\]
All terms have the required weight and form degree.
The first face of each zero-form value is zero.
Its zeroth face is $T^{[w]}$, because
$T^{[N]}T^{[w-N]}=\binom wN T^{[w]}$.
The faces of the one-form values vanish by differentiation.
Thus the displayed assignments preserve both faces and
the coefficient degeneracy, as well as the differential
and all PD operations by the free universal property.
Furthermore the actual identity
$\gamma_N(x)\gamma_{w-N}(x)=\binom wN\gamma_w(x)$
proves $p\psi=p$ on zero-form generators; differentiation
proves it on the other generators. It holds on coefficients
as well.

The pullback therefore has the strict endomorphism
$u(a,c)=(a,C(\psi)(c))$. Its first skeleton is $\psi$
and it satisfies the literal equality $\pi u=\pi$.
The object $T_NG$ and this projection belong to the
ambient admissible category. Inverting $\pi$ there
proves $[u]=1$ in its all-arrow localization.

On the other hand, for every prime $p\geq2N$ the
product term in (S09.154) is zero in (S09.148), so
\[
                        a_p(u)=1+\binom pN>1.
\]
There are infinitely many such primes. Its sequence
is consequently not equal to one outside a finite
set, and $\chi(u)\ne1$ in $\mathcal G$.
The extended functor (S09.153) proves $[u]\ne1$
in the all-arrow localization of $\mathcal E^{[d]}$.
This proves nonfaithfulness by two explicitly specified
morphisms, $u$ and the identity of $E$, and also proves
that this localized groupoid is not terminal.
There is also an explicit diagram killing this very class
inside $\mathcal C^{[d]}$. Let $P_2(N)$ have two copies
$x_w^{(1)},x_w^{(2)}$ of every disk generator of $P(N)$,
with the same endpoints. Evaluate both copies by $p$ and set
\[
 D=T_NG\mathop{\times}_{C((T_NG)_{\leq1})}C(P_2(N)).
\]
The coefficient disk contractions and endpoint surjectivity
are the same ones in A-004. Its full proof, applied to
these two generators per weight, proves that $D$ is admissible
and has the global coefficient comparison through $d$.
Below $N$ every term is literally $S$; above $N$ the same
degree-$d+2$ bounds apply. Thus $D\in\mathcal C^{[d]}$.

The two inclusions of generator copies give actual maps
$i_1,i_2:E\to D$, with the identity on their $T_NG$
coordinate. Folding both copies to $x_w$ gives an actual
map $r:D\to E$. A further map $h:D\to E$ uses
$x_w^{(1)}\mapsto x_w$ and
$x_w^{(2)}\mapsto\psi(x_w)$, again with identity
on $T_NG$. Each assignment commutes with evaluation
because $p\psi=p$, so each is a map of the displayed
pullbacks, preserving every structure. Their exact identities are
\[
             r i_1=r i_2=\mathrm{id}_E,\qquad
             h i_1=\mathrm{id}_E,\qquad h i_2=u.
\]
After inverting $r$, the first pair gives $[i_1]=[i_2]$;
applying $[h]$ gives $[u]=1$ already inside the localization
of $\mathcal C^{[d]}$. This computes the actual loss of
the prime-germ invariant when the doubled first skeleton
is allowed; it asserts no equalization of other parallel maps.
The terms of $E,D$ are countable: their free presentations
have finitely many monomials in each fixed weight, and each
matching diagram and pullback uses finitely many groups.
The same functor restricts to each optional full
subcategory, and all objects and maps just constructed
belong to all of them. This proves every stated variant.
\end{proof}
